\documentclass[10pt]{article}
\usepackage[T1]{fontenc}
\usepackage{lmodern,amsmath,amssymb,amsthm,booktabs,microtype,array}
\usepackage[margin=0.80in]{geometry}
\usepackage[colorlinks=true,urlcolor=blue,linkcolor=blue]{hyperref}
\newtheorem{theorem}{Theorem}
\newtheorem{lemma}{Lemma}
\DeclareMathOperator{\ord}{ord}
\newcommand{\cat}{\mathbin{\Vert}}
\title{A Factorization-Free Certificate for Decimal Mountain Concatenations\\
at Multiplicative Orders at Most 744}
\author{Qichao Wang}
\date{October 18, 2026 \quad\textit{computer-assisted partial result}}
\begin{document}\maketitle
\begin{abstract}
Let \(M(k)=1\cat2\cat\cdots\cat k\cat(k-1)\cat\cdots\cat1\).
We certify that every prime \(p\ne2,5\) with
\(\ord_p(10)\le744\) divides at least one \(M(k)\).
Our finite computation factors neither \(10^r-1\) nor the large
cyclotomic norms appearing in previous research stages.
Instead, we combine two elementary infinite-family arguments
with exact affine-state computations modulo \(10^r-1\),
covering 77 remaining multiplicative orders.
The program checks complete small integer gcd obstructions,
including a nonprimitive-order contaminant at \(r=696\).
This is a bounded-order result, not a solution for all primes.
\end{abstract}
\section{Digit-layer dynamics}
For \(k\ge0\), let \(C_k=1\cat2\cat\cdots\cat k\),
\(R_k=k\cat(k-1)\cat\cdots\cat1\), and let \(D(k)\) be
the total decimal digit count from 1 through \(k\).
For a modulus \(n\) coprime to 10 put
\[
 S_d=(C_{10^d-1},Y_{10^d-1}),\qquad
 Y_k=10^{-D(k)}R_k\pmod n.
\]
Within a full \(d\)-digit layer define
\[
 L=10^{d-1},\quad q=10^d,\quad m=9L,\quad
 A=q^m,\quad B=q-1,\quad
 G=\sum_{i=0}^{m-1}q^i,\quad J=\sum_{i=0}^{m-1}iq^i.
\]
The state transition, valid over any such residue ring, is
\begin{equation}
 (C,Y)\longmapsto
 \bigl(AC+BG-J,\ A^{-1}(Y+LG+J)\bigr).
 \label{eq:layer}
\end{equation}
Only the invertible quantity \(A\) is divided by:
the expression does not presume \(q-1\) is a unit.

For a prime \(p\ge7\), \(p\ne11\), set
\[
 r=\ord_p(10)=u\,v,\quad u=2^a5^c,\quad (v,10)=1,\quad
 T=\operatorname{lcm}(r,\ord_v(10)),\quad\ord_1(10)=1.
\]
Earlier affine-monodromy lemmas imply
\[
 S_{jT}=S_T+(j-1)(U,V),\quad
 (U,V)=S_{2T}-S_T\quad(j\ge1).
 \label{eq:drift}
\]
The normalized mountain residues at \(10^{jT}\) and
\(10^{jT}+1\) are respectively
\(10C_{jT}+Y_{jT}+1\) and \(1000C_{jT}+Y_{jT}+121\).
If \((U,V)\ne(0,0)\), their slopes
\(10U+V\) and \(1000U+V\) cannot both vanish
because \(p\nmid990\).
A linear congruence yields \(1\le j\le p\) and a
divisible mountain at one of the two offsets.

\section{A factorization-free exact-order obstruction}
Fix any integer \(r\ge2\), and set \(N_r=10^r-1\).
The residue class of 10 has exact order \(r\) in
\(\mathbb Z/N_r\mathbb Z\), since \(0<10^h-1<N_r\)
for \(1\le h<r\).
Compute the states in \eqref{eq:layer} in that ring and put
\begin{equation}
 (U_r,V_r)=S_{2T}-S_T,\qquad
 H_r=\gcd(N_r,U_r,V_r).
 \label{eq:gcd}
\end{equation}
For every prime \(p\) of exact order \(r\),
reduction of these integer-ring states modulo \(p\)
gives the prime-field drift. Consequently
\begin{equation}
 (U,V)\equiv(0,0)\pmod p
 \quad\Longrightarrow\quad p\mid H_r.
 \label{eq:obstruction}
\end{equation}

Computing \eqref{eq:gcd} does not require building the huge
\(m=9\cdot10^{d-1}\). Put \(h=r/\gcd(r,d)\) and
write \(m=th+b\), \(0\le b<h\). Since \(q^h=1\),
\[
 G=tG_h+G_b,\qquad
 J=tJ_h+hG_h\binom{t}{2}+thG_b+J_b.
\]
Use modular powering to find
\(m\bmod hN_r\), hence \(t\bmod N_r\) and \(b\).
Binary splitting evaluates \(G_h,J_h,G_b,J_b\)
with \(O(\log r)\) residue-ring operations.
As \(N_r\) is odd, \(\binom t2\bmod N_r\)
is determined by \(t\bmod N_r\).
Thus the entire check uses exact arithmetic without
factoring the repunit \(N_r\).

\begin{theorem}[Certified bounded-order existence]
For every prime \(p\ne2,5\) with
\(\ord_p(10)\le744\), some positive \(k\) satisfies
\(p\mid M(k)\).
For \(p\ge7,\ p\ne11\) we may take
\(k=10^{jT}\) or \(10^{jT}+1\) with \(1\le j\le p\).
\end{theorem}
\begin{proof}
The cases \(p=3\) and \(p=11\) follow respectively
from \(3\mid M(3)=12321\) and \(11\mid M(2)=121\).
For every other \(p\), partition its order \(r\le744\)
according to \(u=2^{v_2(r)}5^{v_5(r)}\) and \(v=r/u\).

If \(u\mid20\), the previously proved elementary
finite-defect transfer gives a perfect-square
first-window polynomial with a subgroup root.
If \(u\nmid20\) and \(v\mid9\), the elementary invariant
\[
 U+V=9\cdot100^{-1}\pmod p\ne0
\]
gives nonzero drift.

Exactly 77 order values \(2\le r\le744\) lie in
neither elementary family. For all 77,
the included programs compute \(H_r\) in
\eqref{eq:gcd} and completely factor this
\emph{small} integer by trial division. They check
\(\ord_q(10)\ne r\) for every prime factor \(q\mid H_r\).
Therefore no prime \(p\) of exact order \(r\)
can satisfy \eqref{eq:obstruction}, and its drift
must be nonzero. The two-offset lemma gives the witness.
\end{proof}

\section{Certificate, audit, and limitations}
The 77 exceptional-to-the-elementary-family orders
consist of 42 previously checked order classes \(r\le500\),
31 new classes \(501\le r\le744\), and four former
high-\(3\)-power cases:
\[
\begin{array}{c|rrrr}
r&216&432&648&675\\\hline
H_r&243&243&729&243
\end{array}
\]
Some representative new values are
\[
\begin{array}{r|rrrrrrrr}
r&504&536&584&624&664&696&725&744\\\hline
H_r&567&99&657&351&9&189&9&27
\end{array}
\]
At \(r=696\), \(H_r=189=3^3\cdot7\);
the extra factor \(7\) is harmless because
\(\ord_7(10)=6\ne696\).
The program therefore tests exact orders of
all gcd prime divisors, rather than assuming the
false stronger assertion that every gcd factor divides
\(990r\).

The 35 additional exact \((r,T,H_r)\) rows are
pinned with SHA-256
\[
\texttt{211b961ab1c7c1762aa39df1e16415a0e56296ef5b4a1a9a9ffbec3d305d5aa3}.
\]
An independent prime-field recurrence verifies
\[
\begin{array}{r|r|r|r}
r&p&U_r\bmod p&V_r\bmod p\\\hline
584&1753&1438&1332\\
664&1993&1889&1956\\
696&13921&13258&9689
\end{array}
\]
The exact, standard-library Python programs are
\texttt{certify\_orders\_le500.py},
\texttt{certify\_orders\_le744.py},
and \texttt{audit\_order\_744\_prime\_field.py},
in the accompanying \texttt{research/} directory.

The bounded-order proof requires independent auditing of
the two previously obtained elementary lemmas and of
the short exact-arithmetic programs.
It neither proves universal nonzero drift
nor resolves OEIS A077187 for all primes.
\end{document}
